\section{Martingalas y bandidos}

Sea \((Y^a_j)_{j\geq1}\) la sucesión i.i.d. de recompensas del brazo \(a\), cuya \(j\)-ésima entrada se revela al elegirlo por \(j\)-ésima vez. Definimos:

\begin{equation} \label{eq:dif_martingala_bandits}
    D^a_t = Y^a_t - \mu_a \quad \forall t\in \mathbbm{N}
\end{equation}

Observemos que es una secuencia de diferencia de martingala \(\forall a\). Sea \(\mathcal{F}^a_t = \sigma(Y^a_1, \cdots, Y^a_t )\).

\begin{proof}
    \begin{align*}
        \E{D^a_{t+1} | \mathcal{F}^a_t} &= \E{Y^a_{t+1} | \mathcal{F}^a_t} - \mu_a \\
                                     &= \E{Y^a_{t+1}} - \mu_a \\
                                    &= \mu_a - \mu_a \\
                                    &= 0 \qedhere.
    \end{align*}
\end{proof}

Luego
\begin{equation} \label{eq:martingala_bandits}
    M^a_t = \sum^t_{s=1}(Y^a_s-\mu_a)
\end{equation}

Es una martingala. Su índice cuenta muestras del brazo, mientras que \(t\) en \(n_t(a)\) cuenta tiradas adicionales. Nos gustaría controlar simultáneamente, para cada brazo \(a\):
\[P\left(\left|\frac{M^a_{n_t(a)}}{n_t(a)}\right| \geq \varepsilon\right).\]


Para eso, introducimos una versión del lema de Hoeffding y el teorema de Azuma-Hoeffding.
